% !TeX root = ../../Book1.tex
\chapter*{序言}
\addcontentsline{toc}{chapter}{序言}
\markboth{序言}{序言}

% 据作者提供的旧线性代数教材第1.1节重新组织；保留数系、数形与方程三条主线。
% 历史引语的来源核对见根目录《序言资料核对.md》。本序兼述五卷安排。

Leopold Kronecker 有一句常被引用的话：
\begin{quote}
上帝创造了整数，其余都是人造的。
\end{quote}

这句话由 Heinrich Weber 的纪念文章《Leopold Kronecker》\cite{Weber1893Kronecker}转述。暂且不论其中的数学哲学立场，它倒很适合作为一次代数学习的开头。我们最早熟悉的数，大抵是用来计数的正整数：有多少个人，有多少本书，数到哪里，便得到一个数。这经验朴素而具体；后来面对的分数、负数、无理数乃至复数，却要逐步建立新的理解。

\ProseParagraphBreak
从正整数 \(1,2,3,\ldots\) 出发，加入表示“没有”的 \(0\)，便得到本书采用的自然数集
\[
\mathbb N=\{0,1,2,3,\ldots\}.
\]
在运算中，零又是加法的单位元：加上零，原来的数不变。这样一个简单性质，后来还会出现在向量、矩阵和多项式的加法中。我们也可以从关于零和后继运算的 Peano 公理出发研究自然数，再在此基础上构造整数、有理数、实数和复数。计算中的常识，由此成为需要说明的定义和命题。

\ProseParagraphBreak
数系扩张是理解代数的一条线索。在自然数中，减法受到限制；引入负数以后，\(a-b\) 不再要求 \(a\ge b\)。整数又不能解出 \(2x=1\)，加入分数，才得到 \(x=1/2\)。但有理数仍不能满足所有求解要求：\(x^2=2\) 没有有理根，\(x^2=-1\) 没有实根。新的问题一次次要求我们扩大原有对象的范围。

\ProseParagraphBreak
这些扩张，各有自己的理由。为使 \(x^2=2\) 有解，添入 \(\sqrt2\) 就已足够；建立整个实数系，还回应了极限与连续性的要求，要补足有理数在这方面的缺陷。复数则让实数之外的平方根成为可以计算的对象。不过，添入一个满足 \(i^2=-1\) 的元素，与证明每个非常数复系数多项式都有复根，是两件不同的事。后者是代数基本定理，不能当成复数的定义。扩张时，原有的加法和乘法应当保持相容；其他性质却未必能够一并保留，例如实数的次序不能延伸成复数上的有序域结构。究竟保留什么、增加什么，须由构造和证明回答。

\ProseParagraphBreak
另一条线索来自数与形的联系。选定坐标以后，方程可以描述几何对象。当 \(a\ne0\) 时，\(ax=b\) 在数轴上确定一个点；在平面上，\(ax+by=c\) 在 \((a,b)\ne(0,0)\) 时确定一条直线。于是线性方程组
\[
\begin{cases}
ax+by=c,\\
dx+ey=f,
\end{cases}
\qquad (a,b)\ne(0,0),\quad(d,e)\ne(0,0),
\]
就是在寻找两条直线的公共点。它们相交、平行或重合，分别对应唯一解、无解或无穷多解。在更多变量中，一个一次项系数不全为零的线性方程描述超平面，消元也就有了寻找这些几何对象公共部分的意义。

\ProseParagraphBreak
Sophie Germain 在《Œuvres philosophiques de Sophie Germain》的《Pensées diverses》中写道\cite[Pensées diverses]{Germain1896}：
\begin{quote}
代数不过是用符号写出的几何，几何不过是用图形表达的代数。
\end{quote}

在线性代数中，这种联系可以直接看见：矩阵记录线性方程，换基改变坐标，而向量空间和线性映射描述坐标背后不变的对象。把一次方程换成多项式方程，便会出现曲线、曲面及更一般的零点集。我们还可以反过来研究这些集合上的多项式函数，考察它们组成的环怎样反映几何性质。本书的交换代数部分，正是沿着这条路，把方程、环和几何对象放在一起研究。

\ProseParagraphBreak
高次方程的求解，则引出第三条线索。二次方程的配方法已经包含一种很有代表性的思想：通过适当的变量变换，使原来混在一起的项变得易于处理。对 \(ax^2+bx+c=0\)、\(a\ne0\)，配方得到
\[
\left(x+\frac{b}{2a}\right)^2=\frac{b^2-4ac}{4a^2},
\qquad
x=\frac{-b\pm\sqrt{b^2-4ac}}{2a}.
\]
当系数为实数而判别式为负时，需要在复数中理解平方根。这里求出的，不只是某一道题的答案，而是一种对整类方程都有效的变形。

\ProseParagraphBreak
三次方程的根式解法通常与 Cardano 的名字相连。Girolamo Cardano 在1545年出版的《Artis magnae, sive de regulis algebraicis liber unus》，通常简称《Ars magna》（《大法》），收录了三次方程的解法；相关发现还包括 Scipione del Ferro 和 Niccolò Tartaglia 的工作，书中也收入了 Lodovico Ferrari 的四次方程解法\cite{Cardano1545}。这些工作所用的语言与今天不同，但其中的变形仍可以用现代记号清楚地看出来。

\ProseParagraphBreak
对一般三次方程 \(ax^3+bx^2+cx+d=0\)，先除以 \(a\ne0\)，再令 \(x=y-b/(3a)\)，可消去二次项，得到 \(y^3+py+q=0\)。若写 \(y=u+v\)，则
\[
y^3+py+q=u^3+v^3+(3uv+p)(u+v)+q.
\]
因此可以要求 \(uv=-p/3\)、\(u^3+v^3=-q\)。这使 \(u^3,v^3\) 成为二次方程 \(T^2+qT-p^3/27=0\) 的两个根。记 \(\Delta=q^2/4+p^3/27\)，便有
\[
u^3=-\frac q2+\sqrt\Delta,\qquad
v^3=-\frac q2-\sqrt\Delta,\qquad uv=-\frac p3.
\]
立方根必须按最后一个等式配对，不能任取两个复立方根再相加。当 \(p\ne0\) 时，可先选 \(u\)，再令 \(v=-p/(3u)\)；此时 \(u\ne0\)，上述两个立方关系相容。取 \(\omega=(-1+i\sqrt3)/2\)，三个根为 \(\omega^k u+\omega^{-k}v\)，其中 \(k=0,1,2\)，重根按重数计入。若 \(p=0\)，直接解 \(y^3=-q\) 即可。最后再把 \(y\) 平移回 \(x\)。

\ProseParagraphBreak
四次方程也可以通过平移消去三次项，化为 \(y^4+py^2+qy+r=0\)。当 \(q=0\) 时，先把它看成关于 \(y^2\) 的二次方程。当 \(q\ne0\) 时，Ferrari 的办法是选择一个适当的 \(m\)，使一次配方后的余项也成为平方。只须令
\[
q^2=4(2m-p)(m^2-r).
\]
这是关于 \(m\) 的三次方程。取其一个复根后，因为 \(q\ne0\)，必有 \(2m-p\ne0\)；在复数中选择 \(s\) 使 \(s^2=2m-p\)，便得到
\[
y^4+py^2+qy+r
=(y^2+m)^2-\left(sy-\frac q{2s}\right)^2.
\]
右边是平方差，因而分解为两个二次式。四次方程的求解就归结为一个三次方程和两个二次方程的求解。这里值得记住的，是如何寻找辅助量，使高次问题化成已经能够处理的低次问题。

\ProseParagraphBreak
到了五次方程，问题发生了变化。Ruffini 和 Abel 的工作所导向的结论是：一般五次及更高次方程不存在由系数出发、只经过有限次四则运算和开方而得到全部根的通用公式。这并不排除某个具体五次方程可以用根式求解。例如，在有理数域上，\(x^5-1=0\) 根式可解，而 \(x^5-x-1=0\) 不根式可解。它们的次数相同，差别须从更细的结构中寻找。这里先把这两个例子作为后续理论要解释的问题；根式可解性的判据安排在本卷\secref{b1:sec:2102}建立。

\ProseParagraphBreak
Galois 理论把注意力转向方程各根之间允许的置换。对于有理系数多项式，保持根之间全部有理系数代数关系的置换组成它的 Galois 群；这个群是否可解，恰好刻画了原方程是否能够用根式求解。由此，求解方程的工作与研究对称性的工作相遇了。为了理解这种联系，我们需要先认识群、正规子群和商群，认识环、理想和域扩张，弄清楚哪些映射保持运算，以及把一个对象分解以后能够保留多少信息。本书从运算与对称讲起，正是要逐步准备这些工具。

\ProseParagraphBreak
我希望这套书能把熟悉的计算与抽象的结构接起来。整数、多项式、矩阵和几何对称，都提供了具体对象；群、环、域、模的定义，则让一些反复出现的规律得到统一表述。抽象的价值，应当在这些对象上看得见：一个同态怎样把问题送到较简单的对象中，一个商怎样把某种差别忽略掉，一个作用又怎样把群的结构变成可以计算的置换或矩阵。读者初遇定义时，不必急着熟悉所有术语，先弄清它所概括的例子，往往更有帮助。

\ProseParagraphBreak
全书拟分五卷展开。第一卷《群、环、域与 Galois 理论》从基本运算、群和对称出发，经过整除、多项式与理想，进入域扩张和 Galois 对应，最终回到根式可解性与尺规作图。开头的群论会先形成一条相对完整的论证路线；读者到达方程理论时，可以重新看见群作用、正规子群和可解性各自发挥的作用。

\ProseParagraphBreak
第二卷《模、多线性代数与非交换环》把线性代数中的标量从域推广到环。向量空间的许多性质需要重新检查，基也不再总是存在。模、张量积和正合列由此成为基本工具，随后讨论半单性、根、Morita 理论和中心单代数。这里建立的模论语言，也为后三卷准备共同基础。

\ProseParagraphBreak
第三卷《交换代数与代数几何方法》以素谱、局部化和有限性为起点，讨论整扩张、维数、完备化、深度和正则性，并用仿射几何解释这些概念。多项式计算、Gröbner 基与消去方法，也在其中占有位置。第四卷《范畴与同调代数》则从对象与映射的组织方式出发，建立复形、消解、\(\operatorname{Ext}\) 与 \(\operatorname{Tor}\)，再进入谱序列、导出范畴和 \(t\)-结构。读到这些内容时，具体的模和映射仍是理解抽象定义的重要依据。

\ProseParagraphBreak
第五卷《表示论：有限群、箭图与李代数》研究怎样用线性变换理解代数对象。从有限群表示与特征标出发，分别发展对称群、箭图和李代数的表示理论，再按需要进入相关专题。整数的计算、方程的对称和向量空间上的作用，在这里又会产生新的联系。五卷所选内容覆盖基础代数与若干主要分支；代数数论、完整的代数几何、代数拓扑及其他相邻学科，各有需要继续学习的体系，本书不以代替这些专门著作为目标。

\ProseParagraphBreak
本书面向已经学习本科线性代数、熟悉集合与映射并接受过基本证明训练的读者。群、环、域仍从定义开始建立。一般点集拓扑不列为全套书的统一先修；无限 Galois 理论与素谱所需的相关知识，会在相应位置准备。各编导读列明直接先修，便于读者按所选路线补齐准备。

\ProseParagraphBreak
五卷的次序，也不是每个人都必须走完的一条长链。初次学习抽象代数，宜先读第一卷的基础部分，再按兴趣继续方程理论，或转入第二卷的模论。第二卷前七章建立的模、正合列、张量积和相关基本工具，可以支撑交换代数、同调代数与表示论的不同路线；进入模论，不要求先读完第一卷的高级域论。第三卷前十四章可以先行阅读，后半涉及深度等内容时，再补第四卷第五至八章的复形、消解及导出函子。第五卷的有限群与李代数核心，也不要求先通读第三、四卷。各编导读会进一步标明直接用到的准备。

\ProseParagraphBreak
带星号的章、节供选读，附录也按选读安排。没有星号，表示它属于本书选定的理论主干，并不意味着适合在第一次阅读时全部完成。一门课程可以从中选取一条连贯路线；自学者则可以先把基础对象和主要定理读熟，再回到较长的证明与专题。阅读建议帮助安排次序，却不能代替证明本身：正文承诺的结论应有完整论证，前瞻和外部调用也应写明其性质。

\ProseParagraphBreak
读一个定义时，可以先写出它的数据：底集合是什么，运算或映射是什么，条件对哪些元素成立。然后找一个最小的例子，再试着去掉某项条件，看原来的构造在哪一步失效。遇到商群、商环或商模，先检查更换代表元是否改变结果；遇到同构，亲自写出映射，查清核与像；遇到泛性质，明确给定了哪些映射、待求映射须满足什么复合等式，再分别检查这样的映射是否存在、是否唯一。这样做，抽象记号才有具体的着落。

\ProseParagraphBreak
读证明，也应当保留计算的习惯。群论中的置换乘法、环论中的整除计算、模论中的矩阵变换，都值得自己动手做几次。一个结论若用到了有限性、交换性或特征条件，就找出这些假设实际出现的那一步。书中的引用应指向具体结果，并说明在当前对象上如何应用；读者不应靠猜测补上关键环节。若仍有一步说不清楚，可以把它单独写成一个待证的小命题，回到定义或所引结论逐项核对。

\ProseParagraphBreak
章末习题是学习的一部分。建议先独立尝试，再阅读所附解答；没有做出时，也先记下自己的构造和受阻之处。读懂解答以后，不妨合上书重新写一遍，尤其留意原先没有想到的对象或映射。例子和反例同样值得积累：以后遇到新定义，手边已有的循环群、置换群、多项式环和有限生成模，常能帮助判断一条猜想是否合理。

\ProseParagraphBreak
伴读资料可以选取 Dummit 与 Foote 的《Abstract Algebra》\cite{DummitFoote2003}和 Lang 的《Algebra》\cite{Lang2002}，再随各分支进入相应专著。同一主题的不同讲法，常在记号、假设和证明顺序上有所区别；比较时应先确认它们讨论的是同一个版本。读过一条完整论证以后，再比较其他讲法，也更容易分辨不同方法的适用范围。

\ProseParagraphBreak
我对读者的期待，是逐渐能够把一个问题说清楚：对象如何定义，运算为何成立，结论依赖什么，构造得到的映射又保留了什么。开始时，一次具体计算可能比一段抽象表述更容易理解；学过一些理论以后，同一段表述又会把许多计算联系起来。希望读者在这两种理解之间往返时，既肯动手核对，也愿意追问理由，并最终能够独立提出和判断新的代数问题。

\bigskip
\begin{flushright}
R. EverStarine

2026年9月
\end{flushright}
